\documentclass[10pt,a4paper]{amsart}
\usepackage[margin=19mm]{geometry}
\usepackage{amsmath,amssymb,mathtools}
\usepackage[scr=rsfs,cal=euler]{mathalfa}
\usepackage{xfrac}
\usepackage[unicode,hidelinks,bookmarks=false]{hyperref}
\newcommand{\ie}{i.e.\ }
\newcommand{\cf}{cf.\ }
\newcommand{\resp}{respectively\ }
\newcommand{\Subsection}{\subsection}
\providecommand{\nobreakdash}{\nobreak\hbox{-}}
\setlength{\emergencystretch}{2em}
\multlinegap0pt
\usepackage{bm}
\usepackage{bbm}
\usepackage{dsfont}
\usepackage{xspace}
\usepackage{cancel}
\usepackage{mathtools}
\usepackage{amsthm}
\makeatletter
\@ifundefined{theorem}{\newtheorem{theorem}{Theorem}}{}
\@ifundefined{lemme}{\newtheorem{lemme}[theorem]{Lemma}}{}
\@ifundefined{proposition}{\newtheorem{proposition}[theorem]{Proposition}}{}
\makeatother
\title[Orthospectrum and simple orthospectrum rigidity: finiteness ]{Orthospectrum and simple orthospectrum rigidity: finiteness and genericity}
\author{Nolwenn Le Quellec}
\date{Source: arXiv:2412.15034v2 (CC BY 4.0)}
\begin{document}
\maketitle

\noindent\emph{Source and licence.} Orthospectrum and simple orthospectrum rigidity: finiteness and genericity. Version arXiv:2412.15034v2 (CC BY 4.0), distributed under the Creative Commons Attribution 4.0 International licence, as recorded in the arXiv licence metadata for this exact version and shown on the abstract page https://arxiv.org/abs/2412.15034v2. Full rights notice: licenses/arxiv-2412.15034-CC-BY-4.0.txt.\par\smallskip

\noindent\textit{Editorial scope.} Counted unit: the Lemme whose source label is \texttt{rel:AxyMin} in the article (source0.tex lines 699--712, source label \texttt{rel:AxyMin}), delivered together with the article's own complete original proof of it (source0.tex lines 714--798, one \texttt{proof} environment of 85 lines). No mathematics is written, repaired or re-derived: statement and proof are the article's own text line for line. Required context retained uncounted: thm:DehnThurston (lines 285--291); in:combiLength (lines 357--363); theorem:Bers (lines 517--523); prop:controleMesureThurston (lines 688--695). Every definition, display and label of those lines is retained unchanged. Omitted from this package and not redistributed: 1--284, 292--356, 364--516, 524--687, 696--698, 713--713, 799--1786. Nothing inside a retained excerpt is omitted or altered: every retained body is delivered exactly as the article's lines are written, so no omission receipt of any kind accompanies this package. Citations: the retained excerpts carry no citation command at all, so no bibliography entry of the article is transcribed and no reference is redistributed. Reference adaptation: none was needed and none was made; every reference inside the delivered unit resolves inside it, so no unresolved reference or citation is produced. Printed slips kept literal: no printed word, symbol, hypothesis or number of the counted statement or of its proof was altered. Layout and class: the article's own class is \texttt{article} with packages including geometry, booktabs, stackengine, fontenc, atkinson, inputenc, bm, graphicx, subcaption, xy, amsthm, amsfonts, mathtools, animate; the wrapper is the lane's pinned \texttt{amsart} editorial wrapper at 10pt on A4, which restates the theorem-like environments the retained text needs and adds the article's own macro definitions that the retained text needs (none), with extra wrapper packages (bm, bbm, dsfont, xspace, cancel, mathtools). The delivered numbering is the unit's own. Assets: no retained span contains a figure environment, a graphics-inclusion command, a PDF-page inclusion, a table or a TikZ picture; no image file is copied into this package, the package declares no asset dependency and no image file is redistributed with it. Licence: version arXiv:2412.15034v2 of this article is distributed under the Creative Commons Attribution 4.0 International licence, as recorded in the arXiv licence metadata for this exact version and shown on the abstract page https://arxiv.org/abs/2412.15034v2. A full rights notice is delivered at licenses/arxiv-2412.15034-CC-BY-4.0.txt. Characters: every retained excerpt is pure ASCII, so the delivered engine, XeLaTeX with its default Latin Modern text font, reports no missing character.
\begin{theorem}\label{thm:DehnThurston}
    For any pair of pants decomposition $\mathcal{P}$ on $S_g^{b}$, the map
    \[
        DT : \mathcal{ML}_{g,b}(\mathbb{Z}) \to Z(\mathcal{P})
    \]
    is a bijection.
\end{theorem}

\begin{theorem}
    Let $\mathcal{P}$ be a pants decomposition on $\chi \in \mathrm{Teich}(S_g^b)$ such that for all $\alpha \in \mathcal{P}$, we have $\ell_\chi(\alpha) \leqslant L$ and $t_\chi(\alpha) \in [0,1]$. Then for any integral multi-curve on $\chi$, we have
    \begin{align}
        \frac{1}{c(L)}\mathcal{L}_\mathcal{P}(\chi,\gamma) \leqslant \ell_\chi(\gamma) \leqslant c(L) \mathcal{L}_\mathcal{P}(\chi,\gamma),\label{in:combiLength}
    \end{align}
    with $c$ that only depends on $L, g$ and $b$.
\end{theorem}

\begin{theorem}\label{theorem:Bers}
    Let $X$ be a finite area hyperbolic surface, possibly with geodesic boundary~$\partial X$. Then $X$ admits a pants decomposition where each curve is of length at most
    \[
        B = \max\{ \ell(\partial X), \mathrm{area}(X) \}.    
    \]
     
\end{theorem}

\begin{proposition}\label{prop:controleMesureThurston}
    For all $g,b \geqslant 0$ and $B \geqslant 2\pi(2g+b-2)$, there exist two constants $c_1(B)$, $c_2(B)>1$ that depend only on $B$ such that for all hyperbolic structure $X$ with geodesic boundary on $S_g^b$ and with $\ell(\partial X) \leqslant B$, we have
    \begin{align*}
        \frac{1}{c_1(B)} \prod_{\gamma : \ell_X(\gamma)\leqslant \varepsilon_{\mathfrak{m}}} \frac{1}{\ell_X(\gamma)S(\ell_X(\gamma))} & \leqslant \mathfrak{m}(X) \mbox{ and,}\\
        \mathfrak{m}(X) \leqslant & c_2(B) \prod_{\gamma : \ell_X(\gamma)\leqslant \varepsilon_{\mathfrak{m}}} \left( 1 + \frac{1}{\ell_X(\gamma)S(\ell_X(\gamma))} + \frac{1}{\min\{\ell_X(\gamma),S(\ell_X(\gamma))\}} \right),
    \end{align*}
    where the products are taken over all simple closed geodesics $\gamma$ of length at most $\varepsilon_{\mathfrak{m}}$ and $S(x)=\sinh^{-1}\left( \frac{1}{\sinh(x/2)} \right)$.
\end{proposition}

\begin{lemme}
    Given $x,y,L>0$, consider the set $A_{x,y}(L)$ defined by
    \[
        A_{x,y}(L)=\{ (m,n) \in \mathbb{Z}_+ \times \mathbb{Z}_+ \mid mx + ny \leqslant L \}.
    \]
    Then for $L>6\max\{x,y\}$, we have
    \begin{align}
        \frac{1}{12}\left( \frac{L^2}{xy} \right) \leqslant \# A_{x,y}(L). \label{rel:AxyMin}
    \end{align}
    And for $L>1$, we have
    \begin{align}
         \# A_{x,y}(L) \leqslant 3L^2\left( 1 + \frac{1}{xy} +\frac{1}{\min\{x,y\}} \right). \label{rel:AxyMaj}
    \end{align}
\end{lemme}

\begin{proof}[Proof of Theorem~\ref{prop:controleMesureThurston}]
    Let $\alpha_1,...,\alpha_s$ be the curves on $X$ of length less than $\varepsilon_{\mathfrak{m}}$. By Theorem~\ref{theorem:Bers}, we can complete this family of curves into a pants decomposition $\mathcal{P}=\{ \alpha_1,...,\alpha_s,...,\alpha_{k} \}$, where $k=3g-3+b$, such that $\ell(\alpha_i) \leqslant B$. Since the Dehn twists on the curves in $\mathcal{P}$ are isometric transformations that change the twist parameters by $1$ we can assume that the twist parameters are between $0$ and $1$. Thus, by inequality~(\ref{in:combiLength}) we have
    
    \begin{multline*}
         \# \{  \gamma \in \mathcal{ML}_{g,b}(\mathbb{Z}) \mid c(B)\mathcal{L}_{\mathcal{P}}(X,\gamma) \leqslant L  \} \\
        \leqslant \# \{ \gamma \in \mathcal{ML}_{g,b}(\mathbb{Z}) \mid \ell_X(\gamma) \leqslant L \} \\ \leqslant \# \{  \gamma \in \mathcal{ML}_{g,b}(\mathbb{Z}) \mid \frac{1}{c(B)}\mathcal{L}_{\mathcal{P}}(X,\gamma) \leqslant L  \}.
    \end{multline*}
    
    Let us start by computing a lower bound for $\mathfrak{m}(X)$.\\
    
    By definition of the combinatorial length and Theorem~\ref{thm:DehnThurston}, we obtain
    {\small
    \begin{align*}
        \# \{  \gamma \in \mathcal{ML}_{g,b}(\mathbb{Z}) \mid c(B)\mathcal{L}_{\mathcal{P}}(X,\gamma) \leqslant L  \} = \# \left\{ (m_i,t_i)_{i=1}^k \in \mathcal{Z}(\mathcal{P})  \mid c(B) \left(  \sum_{i=1}^k m_i S(\ell_X(\alpha_i)) + |t_i|\ell_X(\alpha_i)  \right) \leqslant L  \right\}.
    \end{align*}
    }
    Note that we have $\{ (m_i,t_i)_{i=1}^{k} \mid t_i>0, m_i \in 2\mathbb{Z} \} \subset \mathcal{Z}(\mathcal{P})$, thus
    {\small
    \begin{align*}
        \# \left\{ (m_i,t_i)_{i=1}^k \in (\mathbb{Z}_+ \times \mathbb{Z}_+)^k  \mid c(B) \left(  \sum_{i=1}^k 2 m_i S(\ell_X(\alpha_i)) + t_i\ell_X(\alpha_i)  \right) \leqslant L  \right\}  
        & \leqslant \# \left\{ \gamma \in \mathcal{ML}_{g,b}(\mathbb{Z}) \mid \ell_X(\gamma) \leqslant L \right\} \\
        \# \left\{ (m_i,t_i)_{i=1}^k \in (\mathbb{Z}_+ \times \mathbb{Z}_+)^k \mid c(B) \left( 2 m_i S(\ell_X(\alpha_i)) + t_i\ell_X(\alpha_i)  \right) \leqslant \frac{L}{k}  \right\}  
        & \leqslant  \\
       \prod_{i=1}^{k} \# \left\{ (m_i,t_i) \in  \mathbb{Z}_+ \times \mathbb{Z}_+ \mid   2 m_i S(\ell_X(\alpha_i)) + t_i\ell_X(\alpha_i)   \leqslant \frac{L}{c(B)k}  \right\}  
        & \leqslant  
    \end{align*}
    }
    By Inequality~(\ref{rel:AxyMin}), for $L$ large enough, we have
    
    \begin{align*}
         \prod_{i=1}^{k} \frac{1}{12} \left( \frac{L}{kc(B)} \right)^2 \frac{1}{2S(\ell_X(\alpha_i))\ell_X(\alpha_i)} & 
        \leqslant \# \{ \gamma \in \mathcal{ML}_{g,b}(\mathbb{Z}) \mid \ell_X(\gamma) \leqslant L \} 
    \end{align*}  
        
    We separate the product over the curves $\alpha_i$ into a product over the ones of length at most $\varepsilon_{\mathfrak{m}}$ and a product over the other ones. 
      
    \begin{align*}
         \frac{L^{6g-6+2b}\prod_{i=s+1}^k\frac{1}{S(\varepsilon_{\mathfrak{m}})B}}{(24k^2c(B)^2 )^{3g-3+b}} \prod_{\gamma : \ell_X(\gamma)\leqslant \varepsilon_{\mathfrak{m}}}\frac{1}{\ell_X(\gamma)S(\ell_X(\gamma))} &
        \leqslant \# \{ \gamma \in \mathcal{ML}_{g,b}(\mathbb{Z}) \mid \ell_X(\gamma) \leqslant L \}  \\
         \frac{L^{6g-6+2b}}{(24k^2c(B)^2S(\varepsilon_{\mathfrak{m}})B)^{3g-3+b}} \prod_{\gamma : \ell_X(\gamma)\leqslant \varepsilon_{\mathfrak{m}}}\frac{1}{\ell_X(\gamma)S(\ell_X(\gamma))}  &
        \leqslant 
    \end{align*}  
    We obtain
    \begin{align*}  
         \frac{1}{c_1(B)} \prod_{\gamma : \ell_X(\gamma)\leqslant \varepsilon_{\mathfrak{m}}}\frac{1}{\ell_X(\gamma)S(\ell_X(\gamma))}  &
        \leqslant \frac{\# \{ \gamma \in \mathcal{ML}_{g,b}(\mathbb{Z}) \mid \ell_X(\gamma) \leqslant L \} }{L^{6g-6+2b}}
    \end{align*}
    and letting $L$ go to infinity, we get 
    \[
         \frac{1}{c_1(B)} \prod_{\gamma : \ell_X(\gamma)\leqslant \varepsilon_{\mathfrak{m}}}\frac{1}{\ell_X(\gamma)S(\ell_X(\gamma))}  
        \leqslant \mathfrak{m}(X).
    \]
    Now, let us compute the upper bound on $\mathfrak{m}(X)$. Recall that we have
    \begin{align*}
        \# \{ \gamma \in \mathcal{ML}_{g,b}(\mathbb{Z}) \mid \ell_X(\gamma) \leqslant L \} 
        \leqslant \# \{  \gamma \in \mathcal{ML}_{g,b}(\mathbb{Z}) \mid \frac{1}{c(B)}\mathcal{L}_{\mathcal{P}}(X,\gamma) \leqslant L  \}.
    \end{align*}
    By definition of the combinatorial length and Theorem~\ref{thm:DehnThurston}, we obtain
    {\small
    \begin{align*}
        \# \{ \gamma \in \mathcal{ML}_{g,b}(\mathbb{Z}) \mid \ell_X(\gamma) \leqslant L \} 
        \leqslant & \# \left\{ (m_i,t_i)_{i=1}^k \in \mathcal{Z}(\mathcal{P}) \mid \frac{1}{c(B)}\left(  \sum_{i=1}^k m_i S(\ell_X(\alpha_i)) + |t_i|\ell_X(\alpha_i)  \right) \leqslant L  \right\} \\
          \leqslant & 2^k\# \left\{ (m_i,t_i)_{i=1}^k \in (\mathbb{Z}_+ \times \mathbb{Z}_+)^k \mid \frac{1}{c(B)}\left(  \sum_{i=1}^k m_i S(\ell_X(\alpha_i)) + t_i\ell_X(\alpha_i)  \right)\leqslant L  \right\} \\
          \leqslant & 2^k\# \left\{ (m_i,t_i)_{i=1}^k \in (\mathbb{Z}_+ \times \mathbb{Z}_+)^k \mid \frac{1}{c(B)}\left(  m_i S(\ell_X(\alpha_i)) + t_i\ell_X(\alpha_i)  \right)\leqslant L  \right\} \\
           \leqslant &  \prod_{i=1}^k 2\# \left\{ (m_i,t_i) \in \mathbb{Z}_+ \times \mathbb{Z}_+ \mid  m_i S(\ell_X(\alpha_i)) + t_i\ell_X(\alpha_i)  \leqslant c(B)L  \right\}.
    \end{align*}
    }
    By Inequality~(\ref{rel:AxyMaj}), for $L$ large enough, we have
    {\small
    \begin{align*}
         \# \{ \gamma \in \mathcal{ML}_{g,b}(\mathbb{Z}) \mid \ell_X(\gamma) \leqslant L \} \leqslant &  \prod_{i=1}^k 6(Lc(B))^2\left( 1 + \frac{1}{S(\ell_X(\alpha_i))\ell_X(\alpha_i) } + \frac{1}{\min\{S(\ell_X(\alpha_i)),\ell_X(\alpha_i)\}} \right)\\
          \leqslant & L^{6g-6+2b} \prod_{i=1}^k 6c(B)^2\left(1 + \frac{1}{S(\ell_X(\alpha_i))\ell_X(\alpha_i) } + \frac{1}{\min\{S(\ell_X(\alpha_i)),\ell_X(\alpha_i)\}} \right).
    \end{align*}
    }
    Hence,
    
    \begin{align*}
          \frac{\# \{ \gamma \in \mathcal{ML}_{g,b}(\mathbb{Z}) \mid \ell_X(\gamma) \leqslant L \}}{L^{6g-6+2b}} \leqslant &  c_2(B) \prod_{\gamma : \ell_X(\gamma)\leqslant \varepsilon_{\mathfrak{m}}} \left( 1 + \frac{1}{S(\ell_X(\alpha_i))\ell_X(\alpha_i) } + \frac{1}{\min\{S(\ell_X(\alpha_i)),\ell_X(\alpha_i)\}}  \right).
    \end{align*}
    
    Letting $L$ go to infinity, we obtain 
    \[
      \mathfrak{m}(X) \leqslant c_2(B) \prod_{\gamma : \ell_X(\gamma)\leqslant \varepsilon_{\mathfrak{m}}} \left( 1 + \frac{1}{S(\ell_X(\alpha_i))\ell_X(\alpha_i) } + \frac{1}{\min\{S(\ell_X(\alpha_i)),\ell_X(\alpha_i)\}}  \right).
    \]
\end{proof}

\end{document}
